\documentclass[11pt,a4paper]{article}
\usepackage[margin=2.4cm]{geometry}
\usepackage{booktabs}
\usepackage{graphicx}
\usepackage{xcolor}
\usepackage{hyperref}
\usepackage{enumitem}
\usepackage{titlesec}
\usepackage{fancyhdr}
\usepackage{microtype}
\usepackage{parskip}
\usepackage{fontspec}
\usepackage{polyglossia}
\usepackage{amsmath}
\usepackage{amssymb}

\setmainlanguage{english}
\setotherlanguage{hebrew}
\newfontfamily\hebrewfont{DejaVu Sans}[Script=Hebrew]

\hypersetup{colorlinks=true, linkcolor=blue, urlcolor=blue}

\pagestyle{fancy}
\fancyhf{}
\rhead{Knesset Protocol Journal Project}
\lhead{Step 11 --- Canonicalization Cross-Evaluation}
\rfoot{Page \thepage}
\lfoot{Tomer Morad, Or Israeli, Yoel Marko --- July 2026}

\titleformat{\section}{\large\bfseries}{\thesection.}{0.6em}{}[\titlerule]
\titleformat{\subsection}{\normalsize\bfseries}{\thesubsection}{0.6em}{}

\begin{document}

\begin{center}
  {\LARGE \textbf{Does Canonicalization Help Linking?}}\\[0.3em]
  {\large Step 11 --- a first, imperfect look, on borrowed data}\\[0.6em]
  {\normalsize Tomer Morad \quad Or Israeli \quad Yoel Marko}\\
  {\normalsize July 2026}
\end{center}

\vspace{0.4em}

\section{The question}

Tomer's separate pipeline (\texttt{NLP\_ADV}) rewrites each raw, multi-speaker protocol
segment into a neutral, third-person Hebrew paraphrase --- \textbf{canonicalization} --- using
Gemma-4-31B. The hypothesis: a canonical rewrite should be a cleaner input to an embedding
model than the raw transcript, which is full of speaker turns, procedural boilerplate,
and stylistic noise that has nothing to do with the underlying matter. If true, this should
directly improve Step 10's streaming linking, which lives or dies on embedding quality.

This step is a \textbf{first, quick cross-evaluation} run on whatever data existed at the
time: 19 protocols that Tomer had already segmented and canonicalized independently of our
own gold set. It answers the question provisionally, with two caveats spelled out below,
and sets up Step 12's definitive re-run on our full 523-segment gold benchmark.

\section{Setup}

19 early protocols (Nov--Dec 2022), 656 segments (443 substantive, 213 procedural), each
with both a raw reconstruction (from \texttt{utterance\_ids}) and Tomer's canonical rewrite.
Tomer's LLM-produced \texttt{subject} field is used as a pseudo-topic label for grouping
recurring subjects. Both text variants are embedded with the same
\texttt{multilingual-e5-large} model used throughout Steps 5--10.

\section{Results}

\subsection{Intrinsic topic separation}

Same-subject vs.\ different-subject pair AUC (higher = better-separated embedding space):

\begin{center}
\begin{tabular}{lcc}
\toprule
variant & all pairs & substantive only \\
\midrule
raw & 0.901 & 0.947 \\
\textbf{canonical} & \textbf{0.956} & \textbf{0.986} \\
\bottomrule
\end{tabular}
\end{center}

Canonicalization clearly sharpens topic structure in the embedding space --- substantive-pair
AUC rises from 0.947 to 0.986, the same intrinsic metric used to validate embedding choices
in Steps 3 and 4. Taken alone, this supports Tomer's hypothesis strongly.

\subsection{Downstream streaming linking}

Running our own Step 10 streaming-linking method (recurring subjects as gold) over both
variants:

\begin{center}
\begin{tabular}{lccc}
\toprule
variant & best F1 & precision & recall \\
\midrule
raw & \textbf{0.114} & 0.073 & \textbf{0.261} \\
canonical & 0.095 & \textbf{0.154} & 0.069 \\
\bottomrule
\end{tabular}
\end{center}

Here the picture is \textbf{mixed}. Canonical text roughly doubles precision (0.073
$\rightarrow$ 0.154) --- the right direction under the project's asymmetric cost rule, where
a false merge is worse than a missed link --- but recall collapses, and F1 is lower than raw.
The intrinsic AUC gain does not translate cleanly into a downstream linking win.

\section{Two reasons not to trust this yet}

\begin{enumerate}[leftmargin=1.4em]
\item \textbf{Circularity.} The pseudo-topic \texttt{subject} label used to define
  ``same subject'' is produced by \emph{the same LLM pass} that produces the canonical
  text. Canonical embeddings therefore have a built-in, unearned advantage on any metric
  derived from that label --- the AUC gain in particular should be read as optimistic.
\item \textbf{Scale and label noise.} 19 protocols yield only 29 exact-subject-match
  repeats. Near-duplicate subjects that are really the same underlying matter, but happened
  to get slightly different LLM-generated titles, are counted as \emph{different} --- this
  undercounts true recurrence and makes the linking F1 numbers unstable on this small a
  sample.
\end{enumerate}

\section{What this motivated}

Both caveats point to the same fix: canonicalize \emph{our own} 523 gold segments (Step 7),
where the gold recurrence labels come from independent human annotation, not from the
canonicalizer itself, and where the benchmark has real chains up to length 9 rather than 29
scattered pairs. That re-run --- the clean, definitive version of this experiment --- is
Step 12, and Section~4 of the Step 12 report supersedes the linking numbers above.

\end{document}
